% $Header: /home/vedranm/bitbucket/beamer/solutions/conference-talks/conference-ornate-20min.en.tex,v 90e850259b8b 2007/01/28 20:48:30 tantau $

\documentclass{beamer} %[compress]
% alternatively add handout or draft option

%\mode<handout>%{\setbeamercolor{background canvas}{bg=black!5}}
%\usepackage{pgfpages}
%\pgfpagesuselayout{4 on 1}[a4paper,border shrink=5mm,landscape]
% For handout printing, 4 slides on 1 page

\mode<presentation>
{
  \usetheme{Pittsburgh} % plain template Pittsburgh
  \useoutertheme{infolines} % bottom (author, title, place)
  \useoutertheme[subsection=false]{miniframes} % top (sections, frames)
  \usecolortheme{beaver} %blue color whale
	\setbeamertemplate{frametitle}[default][right] % frame titles on the right
  \setbeamercovered{transparent}
  % or whatever (possibly just delete it)
}

\setbeamersize{text margin left=1cm, text margin right=1cm}
%\setbeamertemplate{navigation symbols}{}
\usepackage[english]{babel}
\usepackage{booktabs}
\usepackage{comment}
% or whatever

\usepackage[cp1250]{inputenc}
% or whatever
\usepackage{subfigure} %subfigures

\usepackage{times}
%\usepackage{breqn}
\usepackage{amsmath}
\usepackage[T1]{fontenc}
\def\sym#1{\ifmmode^{#1}\else\(^{#1}\)\fi} %shortcut for Stata tables
% Or whatever. Note that the encoding and the font should match. If T1
% does not look nice, try deleting the line with the fontenc.

%\logo{\includegraphics[height=0.5cm]{logo.jpg}} % logo on every page



\title[Lecture 5: Basic Tools] % (optional, use only with long paper titles)
{Lecture 5: Basic Meta-Analysis Tools}

\author[T.\ Havranek] % (optional, use only with lots of authors)
{ Tomas Havranek
}
% - Give the names in the same order as the appear in the paper.
% - Use the \inst{?} command only if the authors have different
%   affiliation.

\institute[Prague] % (optional, but mostly needed)
{
  	Charles University, CEPR, METRICS

  %\inst%
  %Economic Research Department\\
  %Czech National Bank
      %\inst{1}%
  %Research Department\\
  %Czech National Bank
  %\and
  %\inst{2}%
  %Institute of Economic Studies\\
  %Charles University, Prague
  
%\pgfdeclareimage[height=1cm]{logo}{logo.jpg} % logo only on title page
%\pgfuseimage{logo}
}% - Use the \inst command only if there are several affiliations.
% - Keep it simple, no one is interested in your street address.

\date[Meta in Econ and Finance] % (optional, should be abbreviation of conference name)
{Research Synthesis in Economics and Finance, \\ University of Canterbury\\[2ex]
  Term 1, 2025
}%\\[4mm]% - Either use conference name or its abbreviation.
% - Not really informative to the audience, more for people (including
%   yourself) who are reading the slides online


%\subject{Intertemporal Substitution}

\begin{comment}
\AtBeginSection[]
{
  \begin{frame}<beamer>{Outline}
    \tableofcontents[currentsection]%,currentsubsection]
  \end{frame}
}
\end{comment}

\begin{document}

\begin{frame}
  \titlepage
\end{frame}

\begin{comment}

\begin{frame}{Outline}
  \tableofcontents
  % You might wish to add the option [pausesections]
\end{frame}

\end{comment}

%\section{Summary}
%\subsection{}

%\section{EIS}
%\subsection{}


\begin{frame}{Mean, median, something else?}
\begin{itemize}
	\item Simple mean of reported estimates: often good starting point
	\item But inefficient, affected by outliers, publication bias
	\item Median: surprisingly robust to most problems in meta-analysis
	\item Classical meta-analysis approach: use inverse variance as the weight
\end{itemize}
\end{frame}


\begin{frame}{More precision -> more weight}
\begin{figure}
	\centering
		\includegraphics[width=1.00\textwidth]{1.pdf}
\end{figure}
\end{frame}


\begin{frame}{More precision -> more weight}
\begin{figure}
	\centering
		\includegraphics[width=1.00\textwidth]{2.pdf}
\end{figure}
\end{frame}


\begin{frame}{More precision -> more weight}
\begin{figure}
	\centering
		\includegraphics[width=1.00\textwidth]{3.pdf}
\end{figure}
\end{frame}


\begin{frame}{Fixed effect}
\begin{figure}
	\centering
		\includegraphics[width=0.9\textwidth]{fixed2.jpg}
\end{figure}
\end{frame}


\begin{frame}{Fixed effect}
    \begin{align*}
        Y_i &= \theta + \varepsilon_i, \tag{1} \\
        V_i &= \frac{\sigma^2}{n}, \tag{2} \\
        W_i &= \frac{1}{V_i}, \tag{3} \\
        M &= \frac{\sum_{i=1}^k W_i Y_i}{\sum_{i=1}^k W_i}, \tag{4} \\
        V_M &= \frac{1}{\sum_{i=1}^k W_i}. \tag{5}
    \end{align*}
		\end{frame}
		
		
		
		
		


% Slide 1: Point Estimate
\begin{frame}{FE vs. UWLS: point estimate}

\begin{itemize}
    \item Both methods yield the same point estimate:
    \[
    \hat{\theta} = \frac{\sum w_i \theta_i}{\sum w_i}, \quad w_i = \frac{1}{\sigma_i^2}
    \]
    \item Key difference: Variance estimation.
\end{itemize}

\bigskip

\begin{table}[]
    \centering
    \begin{tabular}{lc}
        \toprule
        \textbf{Method} & \textbf{Point Estimate} \\
        \midrule
        Fixed-Effect (FE) & \( \hat{\theta} = \frac{\sum w_i \theta_i}{\sum w_i} \) \\
        UWLS & \( \hat{\theta} = \frac{\sum w_i \theta_i}{\sum w_i} \) \\
        \bottomrule
    \end{tabular}
\end{table}

\bigskip

\textbf{Same estimate, different variance!}

\end{frame}

% Slide 2: Variance Comparison
\begin{frame}{FE vs. UWLS: variance}

\renewcommand{\arraystretch}{1.2}  % Increases row spacing for readability
\begin{table}[]
    \centering
    \begin{tabular}{p{3cm} p{3cm} p{3cm}}  % Adjust column widths
        \toprule
        \textbf{Feature} & \textbf{Fixed-Effect (FE)} & \textbf{UWLS} \\
        \midrule
        Assumption  & Identical \(\theta\) & Allows heterogeneity \\
        Variance formula  & \( \frac{1}{\sum w_i} \) & \( \frac{\sum w_i (\theta_i - \hat{\theta})^2}{\sum w_i} \) \\
        Variance components & Within-study only & Within + between-study \\
        Robustness & Sensitive to misspecification & More robust \\
        \bottomrule
    \end{tabular}
\end{table}

\bigskip

\textbf{FE underestimates variance if heterogeneity exists!}

\end{frame}

		


\begin{frame}{Fixed effect}
\begin{figure}
	\centering
		\includegraphics[width=0.9\textwidth]{fixed.jpg}
\end{figure}
\end{frame}


\begin{frame}{Random effects}
\begin{figure}
	\centering
		\includegraphics[width=0.9\textwidth]{random2.jpg}
\end{figure}
\end{frame}


\begin{frame}{Random effects}
    \begin{align*}
        Y_i &= \mu + \xi_i + \varepsilon_i, \tag{1} \\
        V_i &= \frac{\sigma^2}{n}, \tag{2} \\
        W_i &=  \frac{1}{V_i + T^2}, \tag{3} \\
        M &= \frac{\sum_{i=1}^k W_i Y_i}{\sum_{i=1}^k W_i}, \tag{4} \\
        V_M &= \frac{1}{\sum_{i=1}^k W_i}. \tag{5}
    \end{align*}
		\end{frame}
		
		
		
		\begin{frame}{Estimating \(\tau^2\) in random-effects model}

\textbf{Between-study variance} \(\tau^2\) accounts for heterogeneity across studies. Common estimators:
\smallskip
\begin{itemize}
    \item \textbf{DerSimonian-Laird (DL)}:  
    \[
    \tau^2_{DL} = \frac{Q - (k - 1)}{C}, \quad C = \sum w_i - \frac{\sum w_i^2}{\sum w_i}
    \]
    - Simple, but underestimates \(\tau^2\) when \( k \) is small.
    
    \item \textbf{Restricted Maximum Likelihood (REML)}:  
    \[
    \max_{\tau^2} \sum \log(w_i^* + \tau^2) + \sum \frac{(y_i - \hat{\theta})^2}{w_i^* + \tau^2}
    \]
    - More accurate, default in software.
\end{itemize}
\smallskip
\textbf{Rule of thumb:} Use REML unless simplicity is needed!

\end{frame}

		
		
		
\begin{frame}{Random effects}
\begin{figure}
	\centering
		\includegraphics[width=0.9\textwidth]{random.jpg}
\end{figure}
\end{frame}

\begin{frame}{Forest plot}
\begin{figure}
	\centering
		\includegraphics[width=0.9\textwidth]{forest1.png}
\end{figure}
\end{frame}


\begin{frame}{Forest plot}
\begin{figure}
	\centering
		\includegraphics[width=0.45\textwidth]{forest.pdf}
\end{figure}
\end{frame}

\begin{frame}{Box plot}
\begin{figure}
	\centering
		\includegraphics[width=0.45\textwidth]{box1.pdf}
\end{figure}
\end{frame}

\begin{frame}{Box plot}
\begin{figure}
	\centering
		\includegraphics[width=0.65\textwidth]{box2.pdf}
\end{figure}
\end{frame}

\begin{frame}{Box plot}
\begin{figure}
	\centering
		\includegraphics[width=0.95\textwidth]{box3.pdf}
\end{figure}
\end{frame}

\begin{frame}{Histogram}
\begin{figure}
	\centering
		\includegraphics[width=0.9\textwidth]{pattern4.pdf}
\end{figure}
\end{frame}

\begin{frame}{Histogram}
\begin{figure}
	\centering
		\includegraphics[width=0.9\textwidth]{pattern5.pdf}
\end{figure}
\end{frame}

\begin{frame}{Histogram}
\begin{figure}
	\centering
		\includegraphics[width=0.9\textwidth]{pattern6.pdf}
\end{figure}
\end{frame}

\begin{frame}{Cumulative meta-analysis plot}
\begin{figure}
	\centering
		\includegraphics[width=0.9\textwidth]{cumulative.png}
\end{figure}
\end{frame}


\begin{frame}{Funnel plot (topic of the next lecture)}
\begin{figure}
	\centering
		\includegraphics[width=0.9\textwidth]{funnel_pref_all.pdf}
\end{figure}
\end{frame}








\end{document}


